% Include all items presented in the Results section.
%
% The Materials and Methods section should include:
% A clear statement of whether the study is retrospective or prospective. If main prospective study did not include your study purpose in original design, you may classify as a secondary analysis of a prospective trial.
% For prospective studies, enrolled patients become participants. Refer to them as such.
% Date range of the study.
% Sample selection: Describe inclusion and exclusion criteria, consecutive or random selection.
% Indicate how sample size was derived.
% Number, training, and expertise of persons executing and reading the index tests. State the number of years of experience of those who performed readings or evaluations. For resident readings, indicate “in-training.”
% Provide initials for any authors executing index tests or performing readings/evaluations. If those contributors are not authors, mention in acknowledgments but obtain written permission to do so.
% What was evaluated and how was evaluation performed? Mention anonymizing of data or images used in your analysis. Indicate if studies were read independently, if consensus agreements were used, and how disagreements were handled. Single observer studies are discouraged, as are “consensus” reads.
% Description of the pre-specified primary outcome measure(s), as well as secondary and other variables.
%
% For prospective and retrospective studies using human subjects, address in the first paragraph (1) Institutional review board (IRB) approval and patient informed consent. State if consent was obtained; (2) HIPAA compliance for U.S. studies; (3) Clinical Trials Registry Number and Registry Name. Mention whether all or part of your patient population was previously reported with citation of the appropriate reference(s). At the editor’s discretion, we may request a copy/ verification of IRB/ institutional review approval of your study prior to acceptance.
%
% For animal experiments, provide a statement of approval by the institutional animal care committee or appropriate substitute (see the Human and Animal Studies section). If the portion of the animals was previously reported, please indicated the reference.
%
% Describe clearly the number (if animals) and selection of the subjects studied (patients or experimental animals, including controls). Identify any utilized instruments or drugs (including contrast) with trade names, manufacturer’s name in parentheses. Describe procedures in sufficient detail to allow others to reproduce the study. It is essential to explain the manner in which studies were evaluated: independent readings, anonymized or unanonymized to other information, time sequencing between readings of several studies of the same subject to eliminate recall bias, random ordering of studies.
%
% Give references to established methods, including statistical methods that have been published but are not well known; describe new or substantially modified methods and give reasons for using these techniques.
%
% Last Paragraph: State the statistical methods used to analyze the data. Mention the statistical software used, along with version and manufacturer's name. State the p-value used for significance. If statistical analyst is an author, then provide initials. A statement of justification of sample size and/ or power calculation is strongly recommended.
%[radiology]
%[radiology] Radiology recommends that authors seek statistical consultation before planning a study to ensure appropriate enrollment and collection of data and the use of statistical tools. Similarly, discuss studies relating to cost analysis or cost-effectiveness with individuals knowledgeable in these techniques. Radiology reserves the right to have the raw data recalculated by our consultants, as needed. Background information on statistical analyses may be included in the Supplemental Materials when it is essential for understanding the study.
